% Test-case specification TC-0044 (CTest test TC-0044-static).
\documentclass[letterpaper]{article}

\usepackage[letterpaper,margin=1in]{geometry}
\emergencystretch=3em

% Compile from the repository root so this path resolves to
% the pinned WSP submodule.
\input{wsp/testing/test-case-library.tex}

\begin{document}

\def\TCID{TC-0044}
\def\TCTitle{%
  The static library links and answers}
\def\TCPurpose{%
  Verify that a program can link the static form of the library
  (\texttt{wfc\_\allowbreak{}vbrun\_\allowbreak{}static}) instead of loading
  \texttt{wfcvbrun.\allowbreak{}dll}, and that a helper answers.}
\def\TCPriority{\TCPriorityMedium}

\def\TCRequirementRef{%
  REQ-0003}
\def\TCDesignRef{%
  Not applicable.}
\def\TCFeatureRef{%
  \texttt{tests/\allowbreak{}vbrun\_\allowbreak{}static\_\allowbreak{}tests.\allowbreak{}cpp},
  target \texttt{wfc\_\allowbreak{}vbrun\_\allowbreak{}static}}
\def\TCTestTechnique{%
  Requirements-based testing of the link interface.}

\def\TCPreconditions{%
  The library and the test program are built
  (\texttt{scripts/\allowbreak{}Invoke-Build.\allowbreak{}ps1}) and the test
  is registered with CTest.}
\def\TCEnvironment{%
  Windows with the toolchain of
  \texttt{scripts/\allowbreak{}Invoke-Build.\allowbreak{}ps1}; the local
  gate builds and runs x64 Debug, CI builds and runs the other
  architectures. The library is built from this repository in
  \texttt{out/\allowbreak{}build}.}
\def\TCAssumptions{%
  Not applicable.}

\def\TCInputData{%
  The String \texttt{hello} and a null String.}
\def\TCInitialState{%
  No state from earlier runs is needed; the test makes what it needs.}

\def\TCProcedure{%
  \begin{enumerate}
    \item CTest runs \texttt{vbrun\_\allowbreak{}static\_\allowbreak{}tests}.
    \item Link the test with the static library.
    \item Call \texttt{\_\allowbreak{}\_\allowbreak{}vbaLenBstr} and
      \texttt{\_\allowbreak{}\_\allowbreak{}vbaLenBstrB} on \texttt{hello}
      and \texttt{\_\allowbreak{}\_\allowbreak{}vbaLenBstr} on a null
      String.
    \item CTest records the exit code and, on failure, the program's output
      (\texttt{--output-on-failure}).
  \end{enumerate}}
\def\TCExpectedResult{%
  The lengths are 5 characters and 10 bytes, and 0 for the null String; the
  program exits 0 (timeout 120 seconds).}
\def\TCPostConditions{%
  No file or registry state remains that later tests depend on.}

\TCTemplate

\end{document}
